\documentclass[12pt]{article}

\input{../../_assets/preambulo.tex}


\begin{document}

    % 1. Foto de fondo
    % 2. Título
    % 3. Encabezado Izquierdo
    % 4. Color de fondo
    % 5. Coord x del titulo
    % 6. Coord y del titulo
    % 7. Fecha

    
    \input{../../_assets/portada}
    \portadaExamen{ffccA4.jpg}{Análisis\\Matemático II\\Examen I}{AM II. Examen I}{MidnightBlue}{-8}{28}{2026}{Roxana Acedo Parra}

    \begin{description}
        \item[Asignatura] Análisis Matemático II.
        \item[Curso Académico] 2025-26.
        \item[Grado] Doble Grado en Ingeniería Informática y Matemáticas.
        \item[Grupo] Único.
        \item[Profesor] Antonio Cañada Villar.
        \item[Descripción] Primer parcial.
        \item[Fecha] 27 de abril de 2026.
        \item[Duración] 105 minutos.
    \end{description}
    \newpage
    \begin{ejercicio}
        Responde razonadamente a la veracidad o falsedad de las siguientes afirmaciones:
        \begin{enumerate}
            \item (1 punto) Sea $A \subset \mathbb{R}^n$, $\mu(A) > 0$. Entonces existe un intervalo $I = (a, b)$, con $a < b$, tal que $I \subset A$.
            
            \item (1 punto) $\mu(E) = 0 \implies \mathbb{R}^n \setminus E$ es denso en $\mathbb{R}^n$. 
            
            \item (1 punto) Sea $E \subset \mathbb{R}^n$. Si $\mathbb{R}^n \setminus E$ es denso en $\mathbb{R}^n \implies \mu(E) = 0$.
            
            \item (1 punto) Sea $f \in L^1(a, b)$. El conjunto de puntos de $[a, b]$ donde $f$ es discontinua tiene medida cero.
        \end{enumerate}
    \end{ejercicio}

    \begin{ejercicio}[3 puntos]
        Decide, razonadamente, si la función:
        $$ f: (0, 2] \rightarrow \mathbb{R}, \quad f(x) = \frac{\sqrt{x} - \sqrt{x+1}}{x} $$
        pertenece, o no, a $L^1(0, 2)$.
    \end{ejercicio}

    \begin{ejercicio}[3 puntos]
        Decide, razonadamente, si la función:
        $$ f: (2, +\infty) \rightarrow \mathbb{R}, \quad f(x) = \frac{\sqrt{x} - \sqrt{x+1}}{x} $$
        pertenece, o no, a $L^1(2, +\infty)$.
    \end{ejercicio}
    
    \newpage
	\setcounter{ejercicio}{0}

    \begin{ejercicio}
        Responde razonadamente a la veracidad o falsedad de las siguientes afirmaciones:
        \begin{enumerate}
            \item (1 punto) Sea $A \subset \mathbb{R}^n$, $\mu(A) > 0$. Entonces existe un intervalo $I = (a, b)$, con $a < b$, tal que $I \subset A$. \\
            \textbf{Falso.} Que un conjunto tenga medida de Lebesgue positiva $\mu(A) > 0$ no implica que contenga intervalos abiertos. 
            
            \textit{Contraejemplo:} El conjunto de los números irracionales en $[0, 1]$ tiene medida $1$, pero no contiene ningún intervalo $(a, b)$ con $a < b$ porque los racionales son densos en $\mathbb{R}$.
            
            \item (1 punto) $\mu(E) = 0 \implies \mathbb{R}^n \setminus E$ es denso en $\mathbb{R}^n$. \\
            \textit{Nota:} $\mathbb{R}^n \setminus E$ es denso en $\mathbb{R}^n \iff \overline{\mathbb{R}^n \setminus E} = \mathbb{R}^n$ \\
            \textbf{Verdadero.} Si $\mathbb{R}^n \setminus E$ no fuera denso en $\mathbb{R}^n$, existiría una bola abierta $B(x, r)$ tal que $B(x, r) \subset E$. Sin embargo, la medida de cualquier bola abierta es estrictamente positiva ($\mu(B(x, r)) > 0$), lo que contradiría la hipótesis de que $\mu(E) = 0$.
            
            \item (1 punto) Sea $E \subset \mathbb{R}^n$. Si $\mathbb{R}^n \setminus E$ es denso en $\mathbb{R}^n \implies \mu(E) = 0$. \\
            \textbf{Falso.} El hecho de que el complementario sea denso no garantiza que el conjunto original tenga medida nula.
            
            \textit{Contraejemplo:} Sea $E = [0, 1] \setminus \mathbb{Q}$. Su complementario en $[0, 1]$ es $\mathbb{Q} \cap [0, 1]$, que es denso en $[0, 1]$, pero $\mu(E) = 1 \neq 0$.
            
            \item (1 punto) Sea $f \in L^1(a, b)$. El conjunto de puntos de $[a, b]$ donde $f$ es discontinua tiene medida cero. \\
            \textbf{Falso.} Esta propiedad es característica de la integrabilidad de Riemann (Teorema de Lebesgue-Vitali). En la teoría de la medida de Lebesgue, una función puede ser integrable ($f \in L^1$) y ser discontinua en un conjunto de medida positiva o incluso en todo su dominio.
            
            \textit{Ejemplo:} La función característica de los irracionales en $[a, b]$ es integrable Lebesgue, pero discontinua en todos los puntos de $[a, b]$. Definida como 
            $$ f(x) = \begin{cases} 
            1 & \text{si } x \in [a, b] \setminus \mathbb{Q} \\ 
            0 & \text{si } x \in [a, b] \cap \mathbb{Q} 
            \end{cases}$$
            
            Debido a la densidad de los racionales e irracionales, cualquier entorno de cualquier punto $x \in [a, b]$ contiene puntos donde la función vale $1$ y puntos donde vale $0$. Por tanto, $f$ es discontinua en todo su dominio. Al ser discontinua en todos sus puntos, el conjunto de discontinuidades no tiene medida cero ($b-a > 0$). Según el criterio de Lebesgue-Vitali, la función no es integrable Riemann. Sin embargo, según Lebesgue: \\
            
            El conjunto de los racionales $\mathbb{Q} \cap [a, b]$ es numerable y tiene medida de Lebesgue nula ($\mu=0$). Así, $f$ es igual a la función constante $1$ casi por doquier. Su integral es:
            $$ \int_{[a, b]} f \, d\mu = \mu([a, b] \setminus \mathbb{Q}) = (b - a) - 0 = b - a $$
            Como el resultado es finito, $f \in L^1(a, b)$.
        \end{enumerate}
    \end{ejercicio}

    \begin{ejercicio}[3 puntos]
        Decide, razonadamente, si la función:
        $$ f: (0, 2] \rightarrow \mathbb{R}, \quad f(x) = \frac{\sqrt{x} - \sqrt{x+1}}{x} $$
        pertenece, o no, a $L^1(0, 2)$. \\
        
        Como $f$ es continua en $(0,2] \Rightarrow f$ es medible en $(0,2]$.\\
        
        Sabemos que $$L^1(a,b) = \left\{ f:(a, b) \rightarrow \mathbb{R} \text{ medible }: \exists \int^b_a f(x) \, dx\right\}$$ 
        Veamos si acotándola determinamos si existe su integral o no.
        $$ f(x) = \frac{\sqrt{x} - \sqrt{x+1}}{x} = \frac{1}{\sqrt{x}} -\frac{\sqrt{x+1}}{x}$$ 
        De esta manera tenemos que $$\int^2_0 f(x) \, dx = \int^2_0 \frac{1}{\sqrt{x}} \, dx - \int^2_0 \frac{\sqrt{x+1}}{x} \, dx = 2\sqrt{2}-\int^2_0 \frac{\sqrt{x+1}}{x} \, dx$$ 
        Como $\displaystyle \frac{\sqrt{x+1}}{x} \geq \frac{1}{x}$, entonces $$\int^2_0\frac{\sqrt{x+1}}{x} \, dx \geq \int^2_0\frac{1}{x} \, dx = \left[ \ln{|x|}\right]^2_0 = \ln(2) - (-\infty) = +\infty$$
        Consecuentemente, $f \notin L^1(0, 2)$.

        \vspace{1cm}
        
        Aunque \textbf{no sea el objetivo de la asignatura}, se podría hacer calculando la primitiva de la función:
        $$ \int \frac{\sqrt{x} - \sqrt{x+1}}{x} dx = \int \frac{1}{\sqrt{x}} dx - \int \frac{\sqrt{x+1}}{x} dx $$
        Para la segunda integral, usamos el cambio $u = \sqrt{x+1} \implies x = u^2 - 1, dx = 2u \, du$:
        $$ \int \frac{u}{u^2-1} 2u \, du = 2 \int \left( 1 + \frac{1}{u^2-1} \right) du = 2u + \ln\left| \frac{u-1}{u+1} \right| $$
        La primitiva general $F(x)$ es:
        $$ F(x) = 2\sqrt{x} - 2\sqrt{x+1} - \ln\left| \frac{\sqrt{x+1}-1}{\sqrt{x+1}+1} \right| $$
        Analizamos el límite cuando $x \to 0^+$:
        $$ \lim_{x \to 0^+} F(x) = 2(0) - 2(1) - \ln\left| \frac{1-1}{1+1} \right| = -2 - \ln(0) = +\infty $$
        Como la integral diverge en el extremo inferior:
        $$ f \notin L^1(0, 2) $$
    \end{ejercicio}

    \begin{ejercicio}[3 puntos]
        Decide, razonadamente, si la función:
        $$ f: (2, +\infty) \rightarrow \mathbb{R}, \quad f(x) = \frac{\sqrt{x} - \sqrt{x+1}}{x} $$
        pertenece, o no, a $L^1(2, +\infty)$. \\

        Como $f$ es continua en $(2, +\infty) \Rightarrow f$ es medible en $(2, +\infty)$.\\
        
        Podemos volver a acotarla
        $$\frac{\sqrt{x} - \sqrt{x+1}}{x} = \frac{(\sqrt{x} - \sqrt{x+1})\cdot (\sqrt{x} + \sqrt{x+1})}{x\cdot (\sqrt{x} + \sqrt{x+1})}= \frac{-1}{x\cdot (\sqrt{x} + \sqrt{x+1})}$$

        Como $f$ es medible entonces, $f \in L^1(2, +\infty) \iff |f| \in L^1(2, +\infty)$

        $$\frac{1}{x\cdot (\sqrt{x} + \sqrt{x+1})} = \frac{1}{x\cdot\sqrt{x} + x\cdot \sqrt{x+1}} \leq \frac{1}{2x\cdot\sqrt{x}} = \frac{1}{2x^{\frac{3}{2}}}$$

        Calculando $$\int^{+\infty}_2 \frac{1}{2x^{\frac{3}{2}}} \, dx = \left[ \frac{-1}{\sqrt{x}}\right]^{+\infty}_0=0-\left(-\frac{1}{\sqrt{2}}\right) = \frac{1}{\sqrt{2}}$$
        Podemos concluir que $f \in L^1(2, +\infty) $.

        \vspace{1cm}
        
        Aunque \textbf{siga sin ser el objetivo de la asignatura}, podemos evaluar el comportamiento de la primitiva $F(x)$ cuando $x \to +\infty$:
        $$ \int_{2}^{+\infty} f(x) dx = \lim_{x \to +\infty} F(x) - F(2) $$
        Analizamos los términos de $F(x)$ en el infinito: \\
        
        Para $2\sqrt{x} - 2\sqrt{x+1}$:
        $$ \lim_{x \to +\infty} 2(\sqrt{x} - \sqrt{x+1}) = \lim_{x \to +\infty} \frac{-2}{\sqrt{x} + \sqrt{x+1}} = 0 $$
        Para el término logarítmico:
        $$ \lim_{x \to +\infty} \ln\left| \frac{\sqrt{x+1}-1}{\sqrt{x+1}+1} \right| = \ln(1) = 0 $$
        Puesto que $\lim_{x \to +\infty} F(x) = 0$, la integral es igual a $-F(2)$, que es un valor real finito:
        $$ \int_{2}^{+\infty} f(x) dx = - \left( 2\sqrt{2} - 2\sqrt{3} - \ln\left| \frac{\sqrt{3}-1}{\sqrt{3}+1} \right| \right) < \infty $$
        Consecuentemente:
        $$ f \in L^1(2, +\infty) $$
    \end{ejercicio}

    
\end{document}